\section{What, Me Worry?}

\begin{quotex}
The sweet and loving knowledge she says He taught her is mystical theology, that secret knowledge of God which spiritual persons call contemplation. This knowledge is very delightful because it is a knowledge through love. Love is the master of this knowledge and that which makes it wholly agreeable. Since God communicates this knowledge and understanding in the love with which he communicates Himself to the soul, it is very delightful to the intellect, since it is a knowledge belonging to the intellect.
\flright{\textsc{St John of the Cross}, \emph{The Spiritual Canticle}}

What tradition can remain to us of those generations that have not left us a single written line? ... Fortunately, the past never completely dies for man. Man may forget it, but he always preserves it within him. For, take him at any epoch, and he is the product, the epitome, of all the earlier epochs. Let him look into his own soul, and he can find and distinguish these different epochs by what each of them has left behind.
\flright{\textsc{Numa Denis Fustel de Coulanges}, \emph{The Ancient City}}
\end{quotex}


Remember the past ... what if you do? Then what?

\paragraph{Living in the Dark}


\begin{quotex}
Because history---as, moreover, the life of the individual ---is ``worked'' by day and by night. It has a diurnal aspect and a nocturnal aspect. The former is exoteric, whilst the latter is esoteric. The silence and obscurity of the night is always full of events in preparation --- and all that which is unconscious or superconscious in the human being belongs to the domain of ``night''. This is the magical side of history, the side of magical deeds and works acting behind the facade of history ``by day''.
\flright{\textsc{Valentin Tomberg}, \emph{Letter V, The Pope}}

Negative Capability, that is, when a man is capable of being in uncertainties, mysteries, doubts, without any irritable reaching after fact and reason.
\flright{\textsc{John Keats}}
\end{quotex}


There are two kinds of teaching: the speech of day and the knowledge of night, i.e., verbal teaching and direct inspiration. Those who rely only on knowledge based on evidence and reason are afraid of the dark. That is equivalent to looking only at the tip of the iceberg, when there are hidden or occult forces at work.

In John 3, Nicodemus visits Jesus at night. In the night time, the cares of the day should be left behind so the mind can turn to spiritual teachings. Jesus tries to explains the necessity of the second birth to him, but Nicodemus is confused. He can only understand the idea as a daylight teaching rather than as knowledge of the night.

Tomberg explains that there are two truths that must be harmonized:

\begin{itemize}


\item The logic of facts, which deals with the actual, the temporal, the phenomenal

\item Moral logic which is ideal, eternal, and noumenal
\end{itemize}


The discrepancy between the ideal and the actual worlds is a gaping wound, represented by the Sacred Heart. We know what to do to align them, yet we refuse to act.

Calvinists claim the God created sin. However, in Hell, no one blames God for their situation because they understand God's justice. Instead, they spend eternity blaming each other. That is why life on Earth is a foretaste of things to come. We cannot seem to avoid incessant mutual blame. As even Sartre recognized, ``Hell is other people.''

Death is a continuation of life. Hell is more of what you are now.

A woman lost her faith when she learned that her cat would not go to Heaven. Her real concern should be her own postmortem destiny.

\paragraph{The Godforsaken World}


Science displaced the earth from the centre of the universe, but did not remove its godforsakenness. The centre of a circle is the point farthest from the circumference. In that sense the earth is still the centre, and at the centre of the earth, is Hell. Hence, being at the centre is not a privilege, but rather a hazard. It is a long and difficult journey to ascend to God who is even beyond the circumference.

\textbf{St Theresa d'Avila} had a vision of souls falling into Hell like snowflakes falling in a blizzard. It would be hard to argue with her.

The stratification of men has been forgotten. The trend is to classify people by race, whereas the most valid classification is that of caste.

Sin has been forgotten. Yet there are 10 commandments, of quite different characteristics. Atheism and irreligion, for example, are not defensible intellectual positions, but rather moral failings.

\paragraph{The Judge}


\begin{quotex}
Jesus did not trust himself to them because he knew all men and needed no one to bear witness of man; for he himself knew what was in man.
\flright{\textsc{John 2:24-25}}
\end{quotex}


The Real I appears as an impartial Referee or an inner Judge. The referee is just and impartial, as opposed to the personality which is motivated by subjective concerns.

The second coming of Christ takes place in the heart. The third coming appears as the Judge. You can wait until death to be judged, or else you can awaken the judge in this life.

\paragraph{A Serious Life}


At a recent dinner party, time was spent discussing trivia like the best TV shows to watch. On your death bed, I doubt your last words will be, ``I wish I had watched more TV.'' Take your life seriously, do something relevant.

\paragraph{Death Wish}


Sigmund Freud postulated the existence of a Death Instinct, i.e., the desire for unconsciousness. That is unnecessary since the natural state of man is to be unconscious. The real mystery is waking up, or better said, becoming conscious. \textbf{Erich Neumann} makes this observation:

\begin{quotationx}
The unconscious state is the primary and natural one, and the conscious state the product of an effort that uses up libido. There is in the psyche a force of inertia, a kind of psychic gravitation which tends to fall back into the original unconscious condition. In spite of its unconsciousness, however, this state is a state of life and not a state of death. It is just as ridiculous to speak of the death instinct of an apple that falls to the ground as to speak of the death instinct of the ego that falls into unconsciousness.

The pull exerted by the great ``mass'' of the unconscious, i.e., by the collective unconscious with its powerful energy-charge, can only be overcome temporarily by a special performance on the part of the conscious system, though it can be modified and transformed by the building of certain mechanisms.

Even in its waking state our ego consciousness, which in any case forms only a segment of the total psyche, exhibits varying degrees of animation, ranging from reverie, partial attention, and a diffuse wakefulness to partial concentration upon something, intense concentration, and finally moments of general and extreme alertness.
\flright{\textsc{Erich Neumann}, \emph{The Origins and History of Consciousness}}
\end{quotationx}


Modern men falsely believe themselves to be conscious when they are really dominated by subjective elements, fantasies, negative emotions, and so on. In politics, there are two sides: one claiming to be ``woke'', the other to be ``red-pilled''. Both are deluded.

Nevertheless, there is a dim collective consciousness of our ancient ancestors being more conscious. That is the appeal of movies like the \emph{Invasion of the Body Snatchers} and \emph{They Live}.


The first stage of the Hermetic path is to resist the pull into unconsciousness. Hence, there are deliberate exercises to develop attention, concentration, and awareness.

\paragraph{Plebeian Values}


The novel \textit{My Brilliant Friend} is about the friendship of two girls, Elena and Lila, in the slums of Naples after the war. Lila was a child prodigy while Elena was more conventionally intelligent. Nevertheless, there was a telling comment from the school teacher who warned Elena about Lila because Lila was a plebeian.

Each Ancient City had a spiritual foundation, some ancestor or spiritual being who created the constitution for the City. Thus this foundation was essentially spiritual, not genetic. Catholic teaching is that there is an Archangel associated with each city, nation, etc. So the claims of the Ancient City are not implausible; people then had a deeper awareness of transcendent reality.

Each family in the City had its own religious rites; religious observances were an important part of the City's daily life. In addition, each family had clients or dependents. The clients did not have rites of their own but participated in the rites of their family. Eventually, the City would attract other people, alien to the founders.

For example, Athens was on a hill but the plebeians live on the outskirts. The plebeians did not come to Athens to study philosophy, mathematics, or gymnastics for that matter. Rather they were attracted to the security and prosperity of the City. They were the mass without their own religion, unlike the citizens. Ultimately, this led to conflicts as described by Fustel de Coulanges:

\begin{quotationx}
At Rome, the plebeians undertook to form gentes, in imitation of the patricians; at Athens they attempted to overthrow the gentes, to blend them together, and to replace them by the demes, which were established in imitation of them. ... Every gens had a special worship; in Greece the members of the same gens were recognized by the fact that they had performed sacrifices in common from a very early period ... It was a duty to perpetuate this worship from generation to generation, and every man was required to leave sons after him to continue it.

The history of Rome is full of the struggles between the Patricians and Plebeians, struggles that we find in all the Sabine, Latin, and Etruscan cities. We can even remark that the higher we ascend in the history of Greece and Italy, the more profound and the more strongly marked the distinction appears --- a positive proof that the inequality did not grow up with time, but that it existed from the beginning, and that it was contemporary with the birth of cities. What constituted the peculiar character of the plebs was, that they were foreign to the religious organization of the city, and even to that of the family. The plebeian, at first, had no worship, and knew nothing of the sacred family.
\flright{\textsc{Numa Denis Fustel de Coulanges}, \emph{The Ancient City}}
\end{quotationx}


This fact was also noted by Valentin Tomberg. I'm sure he had the French and Russian Revolutions In mind---both hostile to the established religion---when he wrote this:

\begin{quotex}
A pyramid is not complete without its summit; hierarchy does not exist when it is incomplete. Without an Emperor, there will be, sooner or later, no more kings. When there are no kings, there will be, sooner or later, no more nobility. When there is no more nobility, there will be, sooner or later, no more bourgeoisie or peasants. This is how one arrives at the dictatorship of the proletariat, the class hostile to the hierarchical principle, which, however, is the reflection of divine order. This is why the proletariat professes atheism.
\flright{\textsc{Valentin Tomberg}, \emph{Letter IV, The Emperor}}
\end{quotex}


\flrightit{Posted on 2019-04-28 by Cologero}

\begin{center}* * *\end{center}

\begin{footnotesize}
\begin{sffamily}
\texttt{Mike M on 2021-04-28 at 15:51 said:}

``A pyramid is not complete without its summit; hierarchy does not exist when it is incomplete. ... This is why the proletariat professes atheism''

Recently I reflected on my experience with certain existentialist creations, I noted that my experience was entirely different than others in that they were focused on the bleakness and absurdity of it all which would affirm a human striving in the meaninglessness. On the other hand, I took messages of joy, grace, and boundless hope, albeit with some grim realities of the failures of ideologies, pragmatism, and even ``religion''.

I asked a friend why the label existentialist even existed at all, wouldn't any philosophy or work of art have something to do with the fundamentals of human existence? He answered that it begins with atheism. Once atheistic, nihilistic, thought is started, the whole chain of interpretation can only lead to the collapse of hierarchy, reason, structure etc. I saw that the perception of messages and the consequences of those ideas were only too real. Once the Emperor is finished, the whole State collapses as Tomberg notes, my life has not been much different.

In the internal, there is a 1:1 reflection of this process, first goes God and the Emperor, Kings, nobility, working class and guilds, and then all distinction entirely. Then the outcasts and outsiders of oneself have the greatest opportunity to come into one's Personality.

``The first stage of the Hermetic path is to resist the pull into unconsciousness''\\ I then work to stay awake and attentive and discerning what comes out of an unconscious state. Further, accepting responsibility for the actions of the unconscious state instead of attempting to pass the blame onwards. That seems to be a restoration of what is good in the whole system. Perhaps instead of resisting per se (which is necessary at times), its easier when one develops the attraction to consciousness, so to speak.

\hfill

\texttt{Auld Wat on 2022-04-28 at 10:50 said:}

The stratification of man is not hierarchical in the modern sense of the word but rather each caste emanates from the centre to the circumference with the caste closest to the centre representing a singularity or unity. Looking from above we see this more clearly than if we look from the horizontal as moderns do.

Now this may appear to be the reverse of what you say vis a vis Hell but one has to pass through Hell to get to the other side. Dante's journey downward through the circles of Hell is necessary as then, and only then, is he able to begin his ascent up Mount Purgatory to the Celestial City (which itself contains spheres to pass through in turn.

\end{sffamily}
\end{footnotesize}

